\documentclass[10pt,twocolumn]{article}
\usepackage[T1]{fontenc}
\usepackage[utf8]{inputenc}
\usepackage{lmodern}
\usepackage[margin=1.8cm,top=2.4cm,bottom=2cm,columnsep=0.6cm]{geometry}
\usepackage{fancyhdr}
\usepackage{titlesec}
\usepackage{booktabs}
\usepackage{amsmath,amssymb,amsfonts}
\usepackage{microtype}
\usepackage{xcolor}
\usepackage{mdframed}
\usepackage{parskip}
\usepackage{enumitem}
\usepackage{hyperref}

\definecolor{rulered}{RGB}{180,30,30}
\definecolor{boxgray}{RGB}{245,245,245}
\definecolor{keywordgray}{RGB}{100,100,100}

\pagestyle{fancy}
\fancyhf{}
\fancyhead[L]{\textsc{\small Claude Mag}}
\fancyhead[C]{\small\textit{Machine Learning Research Digest}}
\fancyhead[R]{\small 2026-09-07}
\fancyfoot[C]{\small\thepage}
\renewcommand{\headrulewidth}{0.8pt}

\titleformat{\section}{\normalsize\bfseries\color{rulered}}{}{0pt}{}%
  [\vspace{-4pt}\color{rulered}\rule{\columnwidth}{0.4pt}]
\titlespacing{\section}{0pt}{8pt}{4pt}

\hypersetup{colorlinks=true,urlcolor=rulered,linkcolor=black}

\begin{document}
\setlength{\parindent}{0pt}
\setlength{\parskip}{4pt}

{\small\color{keywordgray}\textbf{Keywords:}
  speculative decoding \textbullet{}
  factuality \textbullet{}
  hallucination mitigation \textbullet{}
  inference efficiency}\par\vspace{4pt}

{\LARGE\bfseries\raggedright
  SFAD: Speculative Factuality-Aware Decoding}\par\vspace{4pt}

{\large\itshape\raggedright
  A DPO-Trained Draft Model Plus Epistemic
  Friction Enforces Contextual Faithfulness at
  Speculative Decoding Speeds}\par\vspace{6pt}

{\small\textbf{By} Guanqiao Chen, Di Wang, Lijie Hu
  (MBZUAI et al.)}\par\vspace{2pt}

{\small\color{keywordgray}arXiv:2609.00796 \textbullet{} Preprint, September 2026}
\par\vspace{2pt}\noindent{\color{rulered}\rule{\columnwidth}{0.6pt}}\vspace{6pt}

Contextual faithfulness --- the property that a language model's output
is consistent with a provided source document --- is critical for
retrieval-augmented generation, summarisation, and question answering.
Existing mitigation strategies either double inference cost (contrastive
decoding) or require expensive post-training (RL alignment). SFAD
achieves a 38\% relative reduction in hallucination rate and an 11.6-point
FactScore gain over greedy decoding, at throughput matching standard
speculative decoding, by training a factuality-specialised draft model
with fine-grained preference data and enforcing faithfulness online via
an Epistemic Friction detector.

\begin{mdframed}[backgroundcolor=boxgray,linecolor=rulered,linewidth=1pt,
    innerleftmargin=6pt,innerrightmargin=6pt,
    innertopmargin=6pt,innerbottommargin=6pt]
\textbf{\small AT A GLANCE}\par\vspace{4pt}
\begin{description}[leftmargin=0pt,labelsep=4pt,itemsep=2pt,parsep=0pt,
    font=\small\bfseries]
  \item[Problem:] LLMs hallucinate even when source documents are
    provided; existing remedies sacrifice inference speed.
  \item[Why it matters:] RAG and summarisation systems used in
    production cannot tolerate either hallucinations or latency overhead.
  \item[Method:] DPO-trained context-faithful draft model (ConFide data)
    + Epistemic Friction detector; no change to verifier LLM or workflow.
  \item[Result:] FactScore 72.8\% vs.\ 61.2\% greedy; 88 tok/s
    vs.\ 91 tok/s speculative decoding (no throughput penalty).
  \item[Evidence:] RAGTruth, TruthfulQA, SummEval benchmarks; ablations
    isolating ConFide DPO and Epistemic Friction contributions.
\end{description}
\end{mdframed}

\section{The Big Picture}

Retrieval-augmented generation (RAG) grounds language model outputs in
retrieved source documents. Despite the presence of relevant documents
in context, large language models still hallucinate --- generating
plausible-sounding but contextually unfaithful text that contradicts or
is unsupported by the source.

The challenge is threefold. First, hallucination occurs not because the
model lacks information but because the decoding process does not
explicitly enforce consistency with the context. Second, the most
effective existing methods (contrastive decoding) require two forward
passes, doubling inference cost. Third, post-training alignment via RL
requires significant compute and careful reward design.

SFAD is the first approach to embed factuality enforcement inside
speculative decoding, achieving faithfulness gains with zero throughput
overhead.

\section{Why Existing Approaches Fall Short}

\textbf{Contrastive decoding} compares model outputs under different
prompts (faithful vs.\ non-faithful) and selects tokens with higher
probability under the faithful prompt. It requires two forward passes
per token, doubling inference cost --- impractical in production.

\textbf{RL post-training alignment} trains the model to prefer faithful
outputs via reward models, but requires thousands of feedback examples
and careful reward shaping. The resulting model still hallucinate under
distribution shift.

\textbf{Standard speculative decoding} uses a small draft model to
propose multiple tokens for a large verifier LLM to accept/reject in
parallel. This achieves 2--3$\times$ speedup but makes no claim about
the faithfulness of accepted tokens.

\section{Core Mechanism}

SFAD modifies the draft model in speculative decoding to be
\emph{context-faithful}, so that the token proposals it generates are
more likely to be grounded in the source document. The verifier LLM
then accepts these proposals with high probability --- because they are
both fluent (the draft is a capable LLM) and faithful (the draft is
DPO-trained on faithfulness preferences). No modification to the verifier
or the acceptance criterion is required.

An additional online component, Epistemic Friction, detects tokens
where the draft's proposal diverges from the verifier's preference in
a context-relevant direction, and steers the draft before it proposes ---
further increasing acceptance rate and faithfulness.

\section{Technical Formulation}

\subsection{ConFide Dataset}

ConFide is a preference dataset of (context, chosen output, rejected output)
triples, constructed from model outputs on knowledge-intensive prompts
via \emph{atomic perturbations}:
\begin{itemize}[itemsep=1pt,parsep=0pt]
  \item \textbf{Entity swap}: replace entity names with plausible alternatives
    from the knowledge base.
  \item \textbf{Relation inversion}: change predicate direction
    (e.g.\ ``A was founded by B'' $\to$ ``B was founded by A'').
  \item \textbf{Quantifier change}: replace ``all'' with ``some'',
    exact numbers with nearby values.
  \item \textbf{Temporal shift}: shift dates or tenses.
\end{itemize}
Original outputs (faithful to source) are chosen; perturbed outputs are
rejected. This produces fine-grained atomic contrasts that teach the
draft model exactly which token sequences are unfaithful.

\subsection{DPO Training}

The draft model $\pi_\theta$ is fine-tuned on ConFide with Direct
Preference Optimisation:
\begin{align}
  \mathcal{L}_\mathrm{DPO} = -\mathbb{E}_{(x, y_+, y_-)}\bigl[
  &\log \sigma\bigl(\beta(\log \pi_\theta(y_+|x) - \log \pi_\mathrm{ref}(y_+|x)) \nonumber\\
  &- \beta(\log \pi_\theta(y_-|x) - \log \pi_\mathrm{ref}(y_-|x))\bigr)\bigr]
\end{align}
where $y_+$ is the faithful (chosen) output, $y_-$ is the perturbed
(rejected) output, and $\pi_\mathrm{ref}$ is the pre-DPO draft.

\subsection{Epistemic Friction}

At each proposed token $y_t$, the Epistemic Friction score is:
\[
  \mathrm{EF}(t) =
  \mathrm{KL}\!\left(\pi_\mathrm{draft}(\cdot|x, y_{<t}) \;\|\;
                     \pi_\mathrm{verifier}(\cdot|x, y_{<t})\right)
  \cdot w_\mathrm{specialist}(t)
\]
where $w_\mathrm{specialist}(t)$ is the verifier's mean attention weight
to the source document at position $t$, measuring how context-relevant
the current generation position is.

High EF$(t)$ flags a likely hallucination site. The draft's sampling
temperature is increased at these positions to diversify proposals away
from the unfaithful mode.

\subsection{Speculative Acceptance}

The standard speculative decoding acceptance criterion is unchanged:
\[
  P(\text{accept}\;y_t) = \min\!\left(1,\;
    \frac{\pi_\mathrm{verifier}(y_t|x, y_{<t})}{\pi_\mathrm{draft}(y_t|x, y_{<t})}\right)
\]
Because the DPO draft assigns higher probability to faithful tokens,
the verifier's acceptance rate increases for faithful proposals and
decreases for unfaithful ones --- without any change to the acceptance
formula.

\section{Learning or Inference Procedure}

\textbf{Training:}
\begin{enumerate}[itemsep=2pt,parsep=0pt]
  \item Sample model outputs on knowledge-intensive prompts.
  \item Apply atomic perturbations to construct ConFide preference pairs.
  \item Fine-tune a small draft LLM (3--7B parameters) with DPO on ConFide.
\end{enumerate}

\textbf{Inference:}
\begin{enumerate}[itemsep=2pt,parsep=0pt]
  \item Draft model proposes $k=4$ tokens per speculative step.
  \item Compute EF for each proposed token; re-sample tokens above the
    EF threshold.
  \item Verifier LLM processes the $k$ (steered) tokens in parallel
    and accepts/rejects via the standard criterion.
  \item Accepted tokens are output; rejected tokens restart from the
    last accepted position.
\end{enumerate}

The EF computation adds a single matrix product over the verifier's
cached attention weights --- negligible overhead.

\section{What the Guarantee Says}

The paper proves that under DPO training, the faithful token acceptance
rate improves by at least:
\[
  \Delta_\mathrm{accept} \geq
  \frac{\beta}{\mathcal{Z}} \bigl(\mathbb{E}[\log \pi_\mathrm{DPO}(y_+)]
  - \mathbb{E}[\log \pi_\mathrm{ref}(y_+)]\bigr)
\]
which is strictly positive whenever DPO training increases the faithful
log-probability above the reference model --- a necessary condition for
successful DPO training. The bound is tight for token-level independent
acceptance.

\section{Experimental Findings}

\textbf{Main results across benchmarks:}

\begin{center}
\begin{tabular}{lcc}
\toprule
Method & FactScore$\uparrow$ & Throughput \\
\midrule
Greedy (verifier only)         & 61.2\% & 38 tok/s \\
Contrastive decoding           & 68.4\% & 19 tok/s \\
Standard speculative decoding  & 61.3\% & 91 tok/s \\
\textbf{SFAD (ours)}           & \textbf{72.8\%} & \textbf{88 tok/s} \\
\bottomrule
\end{tabular}
\end{center}

SFAD improves FactScore by 11.6 points over greedy and 4.4 points over
contrastive decoding, while matching standard speculative decoding
throughput. On RAGTruth specifically, hallucination rate drops from
28.3\% to 17.5\% (38\% relative reduction).

\section{Ablations and Interpretation}

\textbf{Without ConFide DPO:} FactScore drops by 6.8 points --- the
draft model fine-tuning is the primary source of faithfulness gain.

\textbf{Without Epistemic Friction:} FactScore drops by 2.1 points,
showing that online steering provides a complementary, additive
faithfulness layer.

\textbf{Atomic vs.\ coarse perturbations in ConFide:} Fine-grained
atomic perturbations outperform sentence-level contrastive pairs by
3.4 FactScore points, validating that granular contrast is essential
for teaching the draft model about specific faithfulness failures.

\textbf{Specialist weight $w_\mathrm{specialist}$:} Attention-based
specialist weighting outperforms uniform weighting by 1.8 FactScore
points. The model is most useful at context-relevant positions,
not uniformly.

\textbf{EF threshold sensitivity:} FactScore is robust to the EF
threshold in the range $[0.3, 0.7]$; below 0.3 too many tokens are
resampled (hurting fluency); above 0.7 too few hallucinations are caught.

\section*{Reference}

Guanqiao Chen, Di Wang, and Lijie Hu.
\textbf{SFAD: Speculative Factuality-Aware Decoding.}
arXiv:2609.00796, September 2026.
\url{https://arxiv.org/abs/2609.00796}

\end{document}
